%! AP Statistics — Unit 8, Lesson 8.5 — Warm-Up (Answer Key)
\documentclass[10pt]{article}
\usepackage{apstats-article}
\usepackage{apstats-key}

\begin{document}

\pageheader{Unit 8, Lesson 8.5}{Warm-Up --- Answer Key}
\namedateperiod

\vspace{0.12in}

A two-way table summarizes data from $n = 180$ observations.

\vspace{0.08in}
\renewcommand{\arraystretch}{1.4}
\begin{tabularx}{0.78\linewidth}{l ccc c}
    & \textbf{Category 1} & \textbf{Category 2} & \textbf{Category 3} & \textbf{Row Total} \\
    \hline
    \textbf{Group A} & 30 & 24 & 6  & 60  \\
    \textbf{Group B} & 50 & 56 & 14 & 120 \\
    \hline
    \textbf{Col Total} & 80 & 80 & 20 & 180 \\
\end{tabularx}

\vspace{0.12in}

\begin{enumerate}

  \item Calculate the expected count for the Group A / Category 1 cell.
        Show your work using the formula $E = \dfrac{(\text{row total})(\text{col total})}{n}$.

        \ansline{$E = (60 \times 80) \div 180 = 4800 \div 180 \approx 26.67$}

  \item Calculate the expected count for the Group B / Category 3 cell.
        Show your work.

        \ansline{$E = (120 \times 20) \div 180 = 2400 \div 180 \approx 13.33$}

  \item Compute all six expected counts and record them in the table below.
        Then check the large-counts condition: is it satisfied?

        \vspace{0.06in}
        \renewcommand{\arraystretch}{1.6}
        \begin{tabularx}{0.78\linewidth}{l ccc}
            & \textbf{Category 1} & \textbf{Category 2} & \textbf{Category 3} \\
            \hline
            \textbf{Group A} & \ans{26.67} & \ans{26.67} & \ans{6.67} \\
            \textbf{Group B} & \ans{53.33} & \ans{53.33} & \ans{13.33} \\
        \end{tabularx}

        \vspace{0.06in}
        Large-counts condition satisfied? \ans{Yes} \quad Explain: \ans{All six expected counts are greater than 5.}

        \begin{teachernote}
            Expected counts: A/C1 = 26.67, A/C2 = 26.67, A/C3 = 6.67;
            B/C1 = 53.33, B/C2 = 53.33, B/C3 = 13.33.
            The smallest is 6.67, which exceeds 5, so the large-counts
            condition is satisfied.
        \end{teachernote}

\end{enumerate}

\end{document}
